\section{Assumptions}
\label{sec:dua_assumptions}
Before delving into the details of the system abstraction layer offered by DUA, it is instrumental to describe the host system configuration that it expects. Recalling what has been said at the beginning of this chapter, the following brief set of assumptions holds mainly for the sake of simplicity: satisfying them allows to employ solutions that are readily available, represent current industry standards in their respective fields, and help to solve significant portions of the problem at hand. What follows from them constitutes a working proof of concept, but nothing prevents the framework from being adapted to different configurations should any requirement change, or should the need arise to support different host system configurations. For now, what is provided in this work satisfies the initial requirements, and any expansion or adaptation is left as future work.

Moreover, for clarity, the following definition is in order here: the term \emph{system configuration}, in this context, refers to a combination of hardware and software that identifies a specific class of systems; each class shall offer and support a defined set of functionalities, hardware components, and operating systems. When a \emph{host} system configuration is mentioned, it refers to a system configuration on which Docker is also present and on which containers are run, \emph{i.e.}, not the one that is found inside containers built and managed by DUA.

\begin{description}
  \item[Hardware platform type] The framework is designed to work on systems that range from single-board computers to more powerful general-purpose computers or even servers, of any architecture, provided that they support the solutions described in \Cref{sec:version_control,sec:containerization}. This assumption excludes lower-end systems like, \emph{e.g.}, microcontrollers, which do play important roles in the field of autonomous systems. This choice has been motivated by two facts: first, microcontrollers and similar systems usually have very limited and constrained sets of resources, which make them incompatible with the software solutions required as of now to build and manage the proposed framework. Second, the variety of microcontroller architectures available today is so vast and each of them is so different from the others, also in terms of operating system support on platforms where this is available, that extending the proposed approach to this class of systems requires a significant effort. This is not to say that the proposed framework can not be adapted to work with microcontrollers, but that is left as future work since their integration requires a dedicated design process.
  \item[Support for version control] The host system should support a version control system, as described in \Cref{sec:version_control}. This is a very important requirement, as the framework relies on version control to manage the software modules that are developed and deployed on the components of the distributed system. The choice of Git as the version control system is motivated by its many features that make it highly efficient when working with projects that embed many modules, as described in \Cref{sec:git}.
  \item[Support for containerization] The host system should support containerization, as described in \Cref{sec:containerization}. This is another very important requirement, as the framework relies on containerization to isolate and manage software components. The choice of Docker as the containerization solution is motivated by its widespread adoption and the functionalities that it offers, enabling it to build replicable environments on a multitude of platforms, as described in \Cref{sec:docker}.
  \item[Operating system support] The framework is designed to natively work on the Linux operating system, as it is the most widely used operating system in the field of distributed autonomous systems. This choice is motivated by the fact that Linux, in its many distributions, offers the widest support for hardware devices and software tools. Moreover, it is built to be developed onto, being natively open-source and offering a stable and fully equipped environment right from the start for the development of software of any kind, especially that which is intended to directly access any hardware resource connected to the host machine. Furthermore, Linux is the operating system that best supports both version control and containerization solutions, whose importance in this project has been highlighted in the previous points. For these same reasons, when a SBC is released to the market, it happens frequently that its natively supported operating system is a Linux distribution, so to offer a consistent environment across all platforms supported by DUA it is natural to assume that the host configurations share this level of similarity. Among the many possible Linux distributions to choose from this work does not impose any specific one, but Ubuntu Linux in one of its latest Long-Term Support (LTS) versions available at the time of writing\footnote{22.04, 24.04.} has been employed for the development and testing of the framework, and is used as the reference for the system abstraction layer described later on. Again, this choice is motivated by similar reasons: Ubuntu is the most popular and easy-to-use Linux distribution, and it is widely supported by many software tools and packages as well as by many SBC vendors. As in any other such case, the use of a different distribution is possible provided that one manages to obtain a working environment with Git, Docker, and any other software dependency that they might require. To conclude, let us note that this assumption is not a strict requirement since the framework can be adapted to work also on Windows and macOS operating systems thanks to Docker, but it is recommended to use a Linux distribution since some features that are crucial for this work are not supported on those OSes due to their different support for containerization which, as explained in \Cref{sec:docker}, always involves a layer of virtualization. Such features are presented and described in the following sections together with their role in the proposed framework but for now, since they mainly concern containerization and map to specific Docker functionalities, a summary table is provided in \Cref{tab:docker_features} in which each feature appears with the name of the corresponding Docker functionality. For completeness, note that while macOS offers a good support to containers running a Linux-based image thanks to Docker Desktop, the best way to run them on Windows is currently to use Docker Desktop in conjunction with the Windows Subsystem for Linux (WSL) \cite{microsoft_windows_nodate}.
  \begin{table}[ht]
    \centering
    \begin{tabular}{|c|c|c|c|}
      \hline
      & \textbf{Linux} & \textbf{Windows} & \textbf{macOS} \\
      \hline
      \textbf{Volume mounts} & \checkmark & \checkmark & \checkmark \\
      \hline
      \textbf{Host networking} & \checkmark & \xmark & \xmark \\
      \hline
      \textbf{Graphical applications} & \checkmark & \checkmark & \xmark \\
      \hline
      \textbf{Hardware devices} & \checkmark & \xmark & \xmark \\
      \hline
      \textbf{Capabilities} & \checkmark & \xmark & \xmark \\
      \hline
      \textbf{IPC namespaces} & \checkmark & \xmark & \xmark \\
      \hline
    \end{tabular}
    \caption{Docker features relevant for the DUA framework and their support on different operating systems.}
    \label{tab:docker_features}
  \end{table}
\end{description}
